\documentclass[11pt]{article}

\usepackage{amsmath,amssymb}
\usepackage[margin=1in]{geometry}
\usepackage{hyperref}

\title{Substitutes and Complements under Product Returns}
\author{Draft notes}
\date{\today}

\begin{document}
\maketitle

\section{Model}

A seller offers two goods separately and as a bundle. The price of good
$i\in\{1,2\}$ is $p_i$. The bundle price is $p_B$, with
$\max(p_1,p_2)\le p_B\le p_1+p_2$. Selling each good alone and selling the
two together is mixed bundling. Producing good $i$ costs $c_i$. Handling a
return of good $i$ costs the seller $h_i$.

There is a unit mass of consumers. Each consumer values good $i$ at $v_i$.
The values are drawn independently from the uniform distribution on $[0,1]$
and are known before buying. Let $W_i$ be the indicator that equals $1$ when
good $i$ works and equals $0$ otherwise. Good $i$ works with probability
$\alpha_i=\mathbb{E}[W_i]$. If good $i$ works, the consumer receives $v_i$.
If it does not work, the consumer receives zero. The values are independent
of $(W_1,W_2)$.

Let $\sigma_i=\sqrt{\alpha_i(1-\alpha_i)}$ be the standard deviation of
$W_i$, and let $\rho$ be the correlation of $W_1$ and $W_2$. The probability
that both goods work is $\alpha_1\alpha_2+\rho\sigma_1\sigma_2$. The
probability that only good $1$ works is
$\alpha_1(1-\alpha_2)-\rho\sigma_1\sigma_2$. The probability that only good
$2$ works is $\alpha_2(1-\alpha_1)-\rho\sigma_1\sigma_2$. The probability
that neither works is $(1-\alpha_1)(1-\alpha_2)+\rho\sigma_1\sigma_2$. These
four probabilities are positive. When $\rho=0$ the goods work independently.
A positive $\rho$ means the goods are complements. They tend to work
together. A negative $\rho$ means the goods are substitutes. One good tends
to work when the other does not.

The consumer buys good $1$, good $2$, the bundle, or nothing, before learning
$(W_1,W_2)$. Returning a good costs the consumer nothing. A consumer who keeps
only good $i$ pays $p_i$. A consumer who keeps both pays $p_B$. A consumer
who keeps nothing pays nothing. The payoff from buying good $i$ alone is
$\max(W_i v_i-p_i,\,0)$. The payoff from buying the bundle is
\[
  \max\bigl(W_1 v_1+W_2 v_2-p_B,\;
  W_1 v_1-p_1,\;
  W_2 v_2-p_2,\;
  0\bigr).
\]
The consumer chooses the option with the highest expected payoff. The
consumer buys nothing when that payoff is negative.

The seller collects the price the consumer pays. The seller pays $c_i$ for
each good $i$ the consumer buys and $h_i$ for each good $i$ the consumer
returns. Profit is expected revenue minus these costs.

\section{Consumer's decision}

The expected payoff from buying nothing is zero. The expected payoff from
buying good $1$ is
\[
  \mathrm{EV}^{1}(v_1,v_2,p_1,p_2)=\alpha_1\max(v_1-p_1,\,0).
\]
The expected payoff from buying good $2$ is
\[
  \mathrm{EV}^{2}(v_1,v_2,p_1,p_2)=\alpha_2\max(v_2-p_2,\,0).
\]

Suppose the consumer buys the bundle and only good $i$ works. Keeping both
pays $v_i-p_B$. This is no larger than $v_i-p_i$, because $p_B\ge p_i$.
The good that does not work has value zero, and returning it pays nothing.
Returning both goods pays zero. The payoff in this state is therefore $\max(v_i-p_i,\,0)$. If
neither good works, the payoff is zero. The expected payoff from the bundle
is
\begin{align*}
  \mathrm{EV}^{b}(v_1,v_2,p_1,p_2)
  ={}&(\alpha_1\alpha_2+\rho\sigma_1\sigma_2)
  \max(v_1+v_2-p_B,\,v_1-p_1,\,v_2-p_2,\,0)\\
  &+(\alpha_1(1-\alpha_2)-\rho\sigma_1\sigma_2)\max(v_1-p_1,\,0)\\
  &+(\alpha_2(1-\alpha_1)-\rho\sigma_1\sigma_2)\max(v_2-p_2,\,0).
\end{align*}

In every state, the bundle payoff is at least the payoff from having bought
only one good. The consumer can keep that good and return the other when it
works and is worth the price, and can return every good otherwise. When both
goods work, keeping both can pay more. Thus
$\mathrm{EV}^{b}(v_1,v_2,p_1,p_2)$ is at least as large as
$\mathrm{EV}^{1}(v_1,v_2,p_1,p_2)$ and at least as large as
$\mathrm{EV}^{2}(v_1,v_2,p_1,p_2)$. The consumer buys the bundle when
$\mathrm{EV}^{b}(v_1,v_2,p_1,p_2)$ is positive. The consumer buys nothing
when that value is zero.

\end{document}
